\section{Visión Artificial y Procesamiento de Imágenes}

\subsection{Modelo de la imagen digital y espacios de color (RGB, HSV) aplicados
al entorno}

Como establecen Gonzalez y Woods \cite{gonzalez2018digital} y Szeliski
\cite{szeliski2022computer}, una imagen digital bidimensional se representa
matemáticamente como una función de intensidad discreta $I(x, y)$, donde $(x, y)$
denotan las coordenadas espaciales en un plano cartesiano y el valor de $I$
corresponde a la amplitud o brillo en dicho punto (píxel). En imágenes multicanal
a color, cada posición espacial está definida por un vector tridimensional en un
espacio vectorial determinado.

El modelo de color predeterminado en los sensores de las cámaras de seguridad es
el espacio RGB (\textit{Red, Green, Blue}). Sin embargo, la representación RGB
acopla fuertemente la información cromática con la intensidad luminosa, haciendo
que variaciones en la luz ambiental alteren drásticamente los tres canales
simultáneamente \cite{gonzalez2018digital}. Para mitigar esta vulnerabilidad en
el parqueadero, el sistema transforma las capturas al espacio de color HSV
(\textit{Hue, Saturation, Value}), el cual desacopla el tono ($H$) y la
saturación ($S$) de la luminancia o brillo ($V$).

Esta separación es fundamental para la robustez del sistema durante la operación
nocturna. Como analizan Liu y Zeng \cite{liu2017vehiclecounting} y Auza Banegas
et al. \cite{auzabanegas2026blpr}, cuando un automóvil aproxima a la barrera de
acceso con los faros delanteros encendidos, se generan destellos directos y
reflexiones de alta intensidad sobre la calzada y la carrocería. Al aislar el
canal de luminancia $V$, el sistema puede aplicar umbralizaciones y análisis
cromáticos sobre las coordenadas de tono y saturación ($H, S$), permitiendo
clasificar el color real del vehículo y localizar la placa de matrícula sin
sufrir saturación ni cegueras provocadas por los haces de luz nocturnos.

\subsection{Filtrado espacial lineal y suavizado Gaussiano}

El filtrado espacial lineal consiste en la aplicación de una operación de
convolución discreta bidimensional entre la matriz de la imagen $I(x, y)$ y un
núcleo o filtro espacial $K(m, n)$ de dimensiones reducidas \cite{szeliski2022computer}:

\begin{equation}
    (I * K)(x, y)
    =
    \sum_{m=-k}^{k} \sum_{n=-k}^{k} I(x - m, y - n) \, K(m, n)
\end{equation}

Entre los filtros de suavizado pasabajo, el \textbf{Filtro Gaussiano} utiliza un
núcleo cuyos coeficientes están definidos por una distribución Gaussiana
bidimensional con desviación estándar $\sigma$. Su función principal es atenuar
los componentes de alta frecuencia espacial mientras preserva las estructuras de
los bordes principales de la escena \cite{gonzalez2018digital}.

En el punto de acceso al parqueadero, los sensores de video capturan
continuamente ruido de alta frecuencia introducido por factores climáticos y
térmicos. La presencia de textura rugosa en el asfalto mojado, el goteo de lluvia
cayendo frente a la lente o el grano fotométrico generado por el sensor en
condiciones de baja iluminación introducen artefactos aleatorios. La aplicación
del suavizado Gaussiano atenúa estos gradientes espaciales espurios,
homogeneizando las superficies antes de ejecutar algoritmos de extracción de
bordes o segmentación de vehículos.

\subsection{Morfología matemática (Top-Hat, Apertura) en el preprocesamiento de
matrículas}

La morfología matemática en procesamiento de imágenes grises se fundamenta en la
teoría de conjuntos y opera mediante la interacción entre la imagen de entrada y
un elemento estructurante $B$ de forma y tamaño definidos
\cite{soille2003morphological}. Las dos operaciones morfológicas no lineales
básicas son la \textit{Erosión} ($\epsilon_B$) y la \textit{Dilatación}
($\delta_B$), a partir de las cuales se construyen la \textit{Apertura}
($\gamma_B = \delta_B \circ \epsilon_B$) —que remueve pequeños brillantes— y el
\textit{Cierre} ($\phi_B = \epsilon_B \circ \delta_B$) —que rellena pequeñas
grietas oscuras—.

Para el realce de estructuras delgadas con iluminación desigual, Soille
\cite{soille2003morphological} formula la transformación \textbf{Top-Hat}
(o Sombrero de Copa):
\begin{itemize}
    \item \textbf{White Top-Hat ($\mathcal{T}_w$):} Diferencia entre la imagen
        original y su apertura ($\mathcal{T}_w(I) = I - \gamma_B(I)$), aislando
        elementos más brillantes que su entorno local.

    \item \textbf{Black Top-Hat / Bottom-Hat ($\mathcal{T}_b$):} Diferencia entre
        el cierre y la imagen original ($\mathcal{T}_b(I) = \phi_B(I) - I$),
        aislando elementos más oscuros que su entorno local.
\end{itemize}

Como destacan Anagnostopoulos et al. \cite{anagnostopoulos2008lpr} y Auza Banegas
et al. \cite{auzabanegas2026blpr}, el estándar de matrículas vehiculares (como el
boliviano) presenta caracteres alfanuméricos oscuros (azules o negros) estampados
sobre un fondo blanco reflectante. Cuando la cámara captura una placa bajo
sombras parciales proyectadas por el parachoques o con iluminación solar
heterogénea, la intensidad del fondo varía drásticamente. La aplicación de la
transformación \textit{Black Top-Hat} aisla y resalta específicamente los trazos
oscuros de las letras y números frente al fondo claro, corrigiendo el gradiente
de iluminación antes de pasar el recorte al motor de Reconocimiento Óptico de
Caracteres (OCR).

\subsection{Limitaciones del procesamiento clásico ante perturbaciones 
ambientales (sombras, oclusión)}

A pesar del rigor matemático del procesamiento digital de imágenes clásico, su
aplicación directa para el control de acceso en entornos no controlados presenta
severas limitaciones estructurales. Como demuestran las investigaciones de
Tatulea et al. \cite{tatulea2019image}, Chandrasekaran et al.
\cite{chandrasekaran2022parking} y Lee et al. \cite{lee2023systematic}, los
métodos deterministas (basados en umbrales estáticos de Otsu, substracción de
fondo o detectores de bordes de Canny) dependen de la suposición rígida de que
las condiciones ópticas de la escena permanecen constantes.

En la realidad operativa de un parqueadero al aire libre, esta premisa se rompe
continuamente por perturbaciones ambientales dinámicas \cite{rioalvarez2019effects}:
\begin{enumerate}
    \item \textbf{Sombras móviles y cambios de iluminación pasante:} El
        movimiento del sol a lo largo del día y la proyección de sombras de
        árboles, estructuras o de los propios vehículos alteran la distribución
        de histogramas de los píxeles, provocando que los algoritmos de
        substracción de fondo identifiquen las sombras como falsos vehículos.

    \item \textbf{Reflejos especulares y charcos en la calzada:} El agua
        acumulada en el pavimento tras una lluvia actúa como un espejo que
        refleja las luces de tráfico y el cielo, generando contornos ruidosos que
        distorsionan las máscaras de segmentación clásica.

    \item \textbf{Oclusiones parciales y variaciones de perspectiva:} Marcos de
        placa, enganches de remolque, acumulación de barro sobre el parachoques
        o ángulos de captura laterales agudos deforman la geometría de los
        objetos. Los métodos basados en plantillas fijas (\textit{template
        matching}) colapsan incapaces de adaptarse a estas deformaciones geométricas.
\end{enumerate}

Estas deficiencias físicas demuestran que las reglas heurísticas ajustadas
manualmente son insuficientes para garantizar la fiabilidad del sistema durante
las 24 horas del día. Esta barrera tecnológica justifica de manera categórica la
necesidad de adoptar modelos de \textbf{Aprendizaje Profundo} (\textit{Deep
Learning}), los cuales reemplazan los filtros diseñados a mano por
representaciones jerárquicas aprendidas directamente desde los datos, logrando
invariancia ante las perturbaciones del mundo real.
